\section{Linear Algebra}
\label{sec:linalg}

Linear algebra is where symbolic computation and fast numerics meet most
directly, and Mathilda meets them with a single idea: a matrix is not a special
kind of value. It is an ordinary expression---a \B{List} of \B{List}s---so the
generic tools you already know (\B{Part}, \B{Map}, \B{Transpose}) apply to it
without ceremony, and a routine like \B{Det} or \B{Inverse} simply reads the
structure off the tree.\index{matrix!as an expression} That uniformity is one of
the two themes of this section. The other is the split that runs underneath every
function here: the \emph{same} call is answered exactly when its entries are exact
and at the speed of C when they are machine numbers. Hand \B{Inverse} an integer
matrix and you get an exact rational inverse; hand it a matrix of measured reals
and the work goes to LAPACK over a packed buffer. The mathematics you type is the
mathematics that runs, whichever world the data lives in.

We climb from structure to spectrum. First the matrix as an expression---its
shape, its constructors, and the packed-array substrate that makes it
fast---then the products and norms that turn vectors into geometry; the
determinant, inverse, and matrix powers; solving linear systems exactly and in
the least-squares sense; rank, row reduction, and the null space, with the
user-facing \texttt{ZeroTest} option that makes exact reduction possible where a
structural pivot would fail; the classical matrix decompositions (LU, QR, SVD,
Schur); and finally the eigenproblem---eigenvalues, eigenvectors, and the Jordan
canonical form---closing with the machine--exact divide that governs all of it.

\subsection{Vectors and matrices as expressions}
\label{ssec:la-matrices}

A vector is a flat list; a matrix is a list of equal-length rows. Because that is
all a matrix is, the everyday expression tools work on it unchanged, and a
handful of predicates and constructors round out the vocabulary.

\begin{codepairs}
\pair{linear-algebra/matrices/1}{A \(3\times3\) matrix is just a list of three
rows.}
\pair{linear-algebra/matrices/2}{\B{Dimensions} reads the shape---here \(3\times3\).}
\pair{linear-algebra/matrices/3}{\B{MatrixQ} confirms it is a rectangular rank-2
array (its siblings \B{VectorQ} and \B{SquareMatrixQ} test the other shapes).}
\pair{linear-algebra/matrices/4}{\B{Part} slices a column: \texttt{m[[All, 2]]} is
the second entry of every row.}
\pair{linear-algebra/matrices/5}{And because rows are ordinary lists, \B{Map}
threads \B{Total} over them to sum each row---no matrix-specific machinery
needed.}
\pair{linear-algebra/matrices/6}{\B{Transpose} exchanges rows and columns.}
\pair{linear-algebra/matrices/7}{\B{ConjugateTranspose} transposes \emph{and}
conjugates---the adjoint \(A^{*}\), the right notion for complex matrices.}
\pair{linear-algebra/matrices/8}{\B{IdentityMatrix} builds the \(n\times n\)
identity.}
\pair{linear-algebra/matrices/9}{\B{DiagonalMatrix} places a list on the
diagonal; the off-diagonal zeros are exact when the diagonal is.}
\pair{linear-algebra/matrices/10}{\B{Diagonal} reads the diagonal back out.}
\end{codepairs}

The lesson of pairs~4 and~5 is worth dwelling on: nothing in \texttt{m[[All, 2]]}
or \texttt{Map[Total, m]} knows it is operating on a ``matrix.'' They are the same
generic operators that act on any expression, and they work here for free because
a matrix \emph{is} an expression. The specialised predicates---\B{SymmetricMatrixQ},
\B{HermitianMatrixQ}, \B{PositiveDefiniteMatrixQ}---layer structural tests on top
of that same representation.

\calloutnote{Under the hood}{Behind the plain list is an invisible fast lane. When
every entry is a machine number, Mathilda may hold the matrix as a
\emph{packed array}\index{packed array}: a flat, row-major buffer of raw
\texttt{double}s carrying its shape on the side, exactly like NumPy's
\texttt{ndarray}, rather than a tree of boxed \texttt{Expr} nodes. \B{Head},
\B{ListQ}, printing, and pattern matching cannot tell a packed list from any
other list---it is the \emph{same value}, only stored densely---so the uniformity
above costs nothing. The visible form of that same storage is \B{NDArray}, which
prints as \texttt{NDArray[\ldots]} and announces itself; the two share every
kernel. Packing and the transparency gate that guards it are covered in
\S\ref{ssec:la-precision}.}

\subsection{Products and norms}
\label{ssec:la-products}

The single operator \B{Dot} (written as an infix period, \texttt{.}) carries the
whole algebra of products: vector--vector, matrix--vector, and matrix--matrix are
one function, contracting the last index of its left argument with the first index
of its right.

\begin{codepairs}
\pair{linear-algebra/products/1}{A matrix times a vector---each output entry is a
row dotted with \(\{x, y\}\).}
\pair{linear-algebra/products/2}{Two matrices multiply to a matrix, row-by-column.}
\pair{linear-algebra/products/3}{Two vectors contract to a scalar, the inner
product \(1\cdot4 + 2\cdot5 + 3\cdot6 = 32\).}
\pair{linear-algebra/products/4}{\B{Cross} is the \(3\)-dimensional cross product:
\(\hat x \times \hat y = \hat z\).}
\pair{linear-algebra/products/5}{In general \B{Cross} takes \(n-1\) vectors of
length \(n\); with one \(2\)-vector it returns the perpendicular in the plane.}
\pair{linear-algebra/products/6}{\B{Tr} sums the diagonal---the trace, invariant
under similarity.}
\end{codepairs}

A \emph{norm} measures length. \B{Norm} defaults to the Euclidean (\(2\)-)norm but
takes any \(p\), and \B{Normalize} scales a vector to unit length.

\begin{codepairs}
\pair{linear-algebra/products/7}{\B{Norm} of a symbolic vector is the Hermitian
\(2\)-norm \(\sqrt{\sum |x_i|^2}\)---note the \B{Abs}, so it is correct for complex
entries.}
\pair{linear-algebra/products/8}{The \(1\)-norm sums absolute values: \(1+2+3=6\).}
\pair{linear-algebra/products/9}{The \(\infty\)-norm takes the largest,
\(\max\{1,2,3\}=3\).}
\pair{linear-algebra/products/10}{\B{Normalize} divides by the norm, sending
\(\{3,4\}\) to the unit vector \(\{3/5, 4/5\}\).}
\end{codepairs}

The \(p\)-norms interpolate a family: \(p=1\) the taxicab length, \(p=2\) the
straight-line length, and \(p=\infty\) the largest coordinate. That the symbolic
answer in pair~7 uses \(\lvert x\rvert^2\) rather than \(x^2\) is not pedantry---for
a complex vector the length-squared is \(\sum x_i\overline{x_i}\), and writing
\(x^2\) would give the wrong answer off the real axis.

\calloutnote{Performance}{On a machine-precision vector or matrix these route
through LAPACK: the \(2\)-norm uses \texttt{dnrm2}'s running-scale recurrence and
the matrix norms use \texttt{dlange}, so a well-scaled vector whose true norm is
representable never overflows to \texttt{inf}. \texttt{Norm[\{1e200, 1e200\}]}
returns \(1.414\times10^{200}\)---the naive \(\sqrt{x_1^2 + x_2^2}\) would square to
infinity first. Both \B{Norm} and \B{Dot} lower inside \B{Compile}, so a norm or
inner product inside a tight numerical loop runs at buffer speed
(\texttt{src/linalg/norm.c}, \texttt{dot.c}).}

\subsection{Determinant, inverse, and matrix powers}
\label{ssec:la-det}

The determinant collapses a square matrix to a single number that vanishes exactly
when the matrix is singular; the inverse undoes the linear map when it does not.

\begin{codepairs}
\pair{linear-algebra/det/1}{\B{Det} of a numeric matrix is exact: \(1\cdot4 -
2\cdot3 = -2\).}
\pair{linear-algebra/det/2}{For a symbolic matrix it returns the Laplace
(cofactor) expansion along the first row.}
\pair{linear-algebra/det/3}{\B{Inverse} of an integer matrix comes back with exact
rational entries.}
\pair{linear-algebra/det/4}{Symbolically, the \(2\times2\) inverse is the classical
adjugate over the determinant \(ad-bc\).}
\pair{linear-algebra/det/5}{Store a matrix to check the defining identity\ldots}
\pair{linear-algebra/det/6}{\ldots{}\(A\,A^{-1}\) is the identity, exactly.}
\pair{linear-algebra/det/7}{\B{MatrixPower} raises a matrix to an integer power by
repeated squaring; the tenth power of \(\left(\begin{smallmatrix}1&1\\1&0\end{smallmatrix}\right)\)
is a block of consecutive Fibonacci numbers.}
\pair{linear-algebra/det/8}{It works symbolically too---here the square of a
general \(2\times2\) matrix.}
\end{codepairs}

Pair~7 is a small gem: the powers of the matrix \(\left(\begin{smallmatrix}1&1\\1&0\end{smallmatrix}\right)\)
generate the Fibonacci sequence, because that matrix encodes the recurrence
\(F_{n+1}=F_n+F_{n-1}\). Its tenth power holds \(F_{11}=89\) and \(F_{10}=55\), and
\B{MatrixPower} reaches them in \(\lceil\log_2 10\rceil=4\) matrix multiplications
rather than nine, by repeated squaring.

\calloutnote{Under the hood}{\B{Det} and \B{Inverse} each choose a path by the kind
of entries they see. An all-integer or rational matrix is handed to FLINT's
\texttt{fmpq\_mat\_det} / \texttt{fmpq\_mat\_inv}, which work in polynomial time and
return the exact answer---so a \(12\times12\) Hilbert determinant is instant, not an
\(O(n!)\) Laplace expansion. A machine-real matrix factorises through LAPACK's
\texttt{dgetrf} LU; a genuinely arbitrary-precision (MPFR) matrix uses an
\(O(n^3)\) MPFR LU. Only a \emph{symbolic} matrix falls to cofactor expansion,
which is why pair~2 shows it. The determinant of a large machine matrix is the
product of the LU pivots and can exceed the IEEE-double range---a
\(200\times200\) random matrix has \(\lvert\det\rvert\approx10^{340}\)---so Mathilda
re-accumulates the pivot product in an MPFR value of unbounded exponent rather than
returning \texttt{inf} (\texttt{src/linalg/det.c}, \texttt{inv.c}, \texttt{matpow.c}).}

\subsection{Solving linear systems}
\label{ssec:la-linsolve}

To solve \(A\mathbf{x}=\mathbf{b}\) is the central task of the subject. \B{LinearSolve}
does it exactly when it can, over- and under-determined systems included; when no
exact solution exists, \B{LeastSquares} returns the closest thing.

\begin{codepairs}
\pair{linear-algebra/linsolve/1}{A square system, solved exactly over the
rationals.}
\pair{linear-algebra/linsolve/2}{With symbolic entries the answer is Cramer's
rule: each component is a ratio of determinants sharing the denominator
\(ru-st=\det A\).}
\pair{linear-algebra/linsolve/3}{An \emph{under}-determined system (more unknowns
than equations) returns one particular solution, with the free variables set to
zero---use \B{Solve} for the full parametrised family.}
\pair{linear-algebra/linsolve/4}{When the system is \emph{over}-determined and
inconsistent, \B{LeastSquares} minimises \(\lVert A\mathbf{x}-\mathbf{b}\rVert\):
the best-fit line \(y = \tfrac{19}{3} + \tfrac12 x\) through \((1,7),(2,7),(3,8)\).}
\pair{linear-algebra/linsolve/5}{And when the system \emph{is} consistent, the
least-squares solution coincides with the exact one.}
\end{codepairs}

The distinction in pairs~3 and~4 is the whole geometry of the problem.
\B{LinearSolve} answers ``is \(\mathbf{b}\) in the column
space\index{Cramer's rule} of \(A\)?''---if so,
which combination of columns produces it. When \(A\) has more columns than rows the
answer is a whole affine subspace, and \B{LinearSolve} hands back one point of it;
when \(A\) has more rows than columns and \(\mathbf{b}\) lies off the column space,
there is no solution at all, and \B{LeastSquares} instead returns the
\(\mathbf{x}\) whose image \(A\mathbf{x}\) is the orthogonal projection of
\(\mathbf{b}\) onto the column space---the least-squares fit\index{least squares} that underlies all of
linear regression (and \B{Fit}, in \S\ref{sec:statistics}).

\calloutnote{Theory}{\B{LeastSquares} minimises \(\lVert A\mathbf{x}-\mathbf{b}\rVert_2\).
Setting the gradient to zero gives the \emph{normal equations} \(A^{*}A\,\mathbf{x} =
A^{*}\mathbf{b}\), whose solution is \(\mathbf{x}=A^{+}\mathbf{b}\) with \(A^{+}\) the
\B{PseudoInverse}\index{Moore-Penrose pseudoinverse}. For a rank-deficient \(A\)
the minimiser is not unique, and \(A^{+}\mathbf{b}\) picks out the one of smallest
norm. A machine-real system takes a thin SVD rather than forming \(A^{+}\)
explicitly---a \(500\times500\) least-squares fit that never finished on the exact
pipeline costs \(77\)~ms this way (\texttt{src/linalg/lstsq.c}).}

\usagebox{LinearSolve}

\subsection{Rank, row reduction, and the null space}
\label{ssec:la-reduction}

Beneath solving lies elimination. \B{RowReduce} brings a matrix to reduced row
echelon form; \B{MatrixRank} counts the independent rows; and \B{NullSpace}
returns a basis for the vectors the matrix annihilates. The three are one
computation seen from three angles.

\begin{codepairs}
\pair{linear-algebra/reduction/1}{\B{RowReduce} of a rank-deficient matrix: the
bottom row of zeros is the dependency made visible.}
\pair{linear-algebra/reduction/2}{\B{MatrixRank} reports \(2\)---the third row is
the sum of a scaling of the first two.}
\pair{linear-algebra/reduction/3}{\B{NullSpace} returns the one vector \(\{1,-2,1\}\)
spanning the kernel.}
\pair{linear-algebra/reduction/4}{Store the matrix\ldots}
\pair{linear-algebra/reduction/5}{\ldots and confirm it: \(A\) times the null
vector is \(\mathbf 0\).}
\pair{linear-algebra/reduction/6}{These are exact even symbolically---row-reducing
a matrix whose third row is \(\text{row}_1+\text{row}_2\) exposes the zero row and
the exact rational entries above.}
\pair{linear-algebra/reduction/7}{\B{MatrixRank} sees the symbolic dependency too:
the second row is \(2\) times the first, so the rank is \(1\).}
\end{codepairs}

The rank-nullity theorem\index{rank-nullity theorem} is on display: a
\(3\times3\) matrix of rank \(2\)
(pair~2) has a null space of dimension \(3-2=1\) (pair~3), and their sum recovers
the column count. That Mathilda finds these exactly, over rationals \emph{and} over
symbols, is the point of pairs~6 and~7---the elimination is fraction-free, so no
spurious denominator ever appears.

\calloutnote{Under the hood}{The default reduction is a Bareiss-like
\emph{fraction-free} Gauss--Jordan elimination\index{Bareiss algorithm}: it clears
each column with cross-multiplications chosen so that the exact divisions leave no
denominator larger than necessary, keeping intermediate coefficients from
exploding the way naive Euclidean elimination would. An all-integer or rational
matrix instead goes to FLINT's \texttt{fmpq\_mat\_rref}, whose reduced echelon form
is unique and therefore matches. \B{MatrixRank} and \B{NullSpace} both call back
into \B{RowReduce} through the evaluator, so every improvement to it propagates to
them for free (\texttt{src/linalg/linsolve.c}, \texttt{matrank.c},
\texttt{nullspace.c}).}

\subsection{Exact reduction and the \texttt{ZeroTest} option}
\label{ssec:la-zerotest}

Elimination lives or dies on one question asked at every pivot: \emph{is this entry
zero?} For numbers it is trivial. For symbols it is the hardest question in the
system, and it is where exact linear algebra quietly goes wrong if you are not
careful. Consider a matrix whose two rows are equal---but equal only after a
trigonometric identity that Mathilda does not apply on its own.

\begin{codepairs}
\pair{linear-algebra/zerotest/1}{The Pythagorean identity is true, but it does not
auto-simplify: Mathilda leaves \texttt{Cos[x]\^{}2 + Sin[x]\^{}2} as written.}
\pair{linear-algebra/zerotest/2}{So these two rows are mathematically identical---the
top-left entry \emph{is} \(1\)---but not \emph{structurally} identical.}
\pair{linear-algebra/zerotest/3}{The default \B{NullSpace} uses a structural zero
test. It cannot see that the rows are dependent, so it wrongly reports full rank
and an empty kernel.}
\pair{linear-algebra/zerotest/4}{Supply \texttt{ZeroTest -> (Simplify[\#] === 0 \&)}
and the reduction consults \B{Simplify} at every pivot---now it finds the
dependency and returns the true null vector.}
\end{codepairs}

This is the reason for the user-facing \texttt{ZeroTest} option on \B{RowReduce}
and \B{NullSpace}.\index{ZeroTest option}\index{zero testing!symbolic} The default
structural test treats each distinct radical or transcendental as an independent
symbol---fast, and correct whenever the answer is structurally apparent, but blind
to any zero that only appears after simplification. Passing a predicate---a body
like \texttt{(Simplify[\#] === 0 \&)} or a head like \B{PossibleZeroQ}---replaces
that test with a semantic one, and the elimination pivots on genuine, not merely
apparent, non-zeros. There is no \emph{perfect} zero test---that this must be so is a
theorem, not a shortfall of effort---so how hard to test at each pivot is a choice left to
you; \S\ref{ssec:alg-possiblezeroq} takes up \mcode{PossibleZeroQ} and the undecidability
theorem behind it.

The same fault line runs through \emph{numeric} reduction, where the question
becomes ``is this entry small enough to treat as zero?'' Here the answer depends on
a tolerance, and the exact and machine paths part company sharply.

\begin{codepairs}
\pair{linear-algebra/zerotest/5}{A matrix whose two off-diagonal pivots are tiny
but exactly non-zero.}
\pair{linear-algebra/zerotest/6}{Exactly, this matrix has three independent
rows---the tiny entries \(10^{-10}\) and \(10^{-20}\) are genuinely non-zero.}
\pair{linear-algebra/zerotest/7}{Convert to machine reals and one pivot falls below
the default tolerance: \B{MatrixRank} now reports \(2\).}
\pair{linear-algebra/zerotest/8}{Force \texttt{Tolerance -> 0} and every non-zero
counts again, recovering the rank of \(3\).}
\end{codepairs}

Neither answer is wrong; they answer different questions. Exact reduction asks
whether an entry is \emph{algebraically} zero and keeps the tiny pivots; machine
reduction asks whether it is zero \emph{to working precision} and discards them,
which is exactly what you want when the small entries are round-off rather than
signal. The lesson is to know which question you are asking---and to reach for
\texttt{Tolerance} (numeric) or \texttt{ZeroTest} (exact) when the default answers
a different one.

\subsection{Matrix decompositions}
\label{ssec:la-decomp}

A decomposition rewrites a matrix as a product of simpler factors---triangular,
orthogonal, diagonal---and almost every serious numerical algorithm is a
decomposition in disguise. Mathilda ships the classical four.

\begin{codepairs}
\pair{linear-algebra/decomp/1}{\B{LUDecomposition} returns the combined
lower/upper factors, the pivot permutation, and a condition estimate---here the
factorisation of a Vandermonde-like matrix.}
\pair{linear-algebra/decomp/2}{\B{QRDecomposition} factors \(A\) into an
orthonormal \texttt{q} and an upper-triangular \texttt{r}. This is the classic
textbook \(3\times3\), and its \(Q\) is exactly rational.}
\pair{linear-algebra/decomp/3}{The orthonormal factor \texttt{q}: its rows are
mutually perpendicular unit vectors.}
\pair{linear-algebra/decomp/4}{The upper-triangular factor \texttt{r}.}
\pair{linear-algebra/decomp/5}{And they reconstruct the original exactly:
\(Q^{\mathsf T} R = A\).}
\pair{linear-algebra/decomp/6}{\B{SingularValueDecomposition} of a rank-\(1\)
matrix: a single non-zero singular value \(\sqrt{10}\) on the diagonal of
\(\Sigma\), the rest zero.}
\pair{linear-algebra/decomp/7}{A symmetric machine matrix for the Schur form\ldots}
\pair{linear-algebra/decomp/8}{\ldots factored by \B{SchurDecomposition}\ldots}
\pair{linear-algebra/decomp/9}{\ldots into an orthonormal \texttt{q} and an
upper-triangular \texttt{t} (diagonal here, since the matrix is symmetric): its
diagonal carries the eigenvalues.}
\pair{linear-algebra/decomp/10}{\B{Chop} confirms \(A = Q\,T\,Q^{*}\) to machine
zero.}
\end{codepairs}

Each factorisation answers a different need. \B{LUDecomposition} is the engine of
\B{LinearSolve}: once \(A=LU\), each right-hand side costs only two triangular
solves. \B{QRDecomposition} orthogonalises the columns and is the stable way to
solve least-squares problems and the workhorse of the eigenvalue algorithm.
\B{SingularValueDecomposition} is the most revealing of all---the singular values
are the semi-axes of the ellipsoid \(A\) maps the unit sphere onto, the number of
non-zero ones is the rank (pair~6 makes a rank-\(1\) matrix show a single
\(\sqrt{10}\)), and truncating the small ones is the mathematics of data
compression and \B{PseudoInverse}. \B{SchurDecomposition}\index{Schur decomposition} is the theoretical keystone:
\emph{every} square matrix is unitarily
similar to a (quasi-)upper-triangular one whose diagonal is the spectrum, so the
eigenvalues can be read off without ever solving a polynomial.

\calloutnote{Under the hood}{The decompositions dispatch on precision. A
machine-real matrix goes to LAPACK---\texttt{dgetrf} (LU), a Householder
\texttt{dgeqrf} (QR)\index{Householder reflection}, divide-and-conquer
\texttt{dgesdd} (SVD), and the Francis QR iteration \texttt{dgees} (Schur)---wired
through a four-tier autodetection ladder (Apple Accelerate, then two flavours of
system LAPACK, then a graceful fallback). An arbitrary-precision matrix runs a
hand-rolled MPFR kernel at the requested precision (Householder QR, one-sided
Jacobi SVD, Hessenberg+Francis Schur). An exact or symbolic matrix stays on a
modified Gram--Schmidt or eigendecomposition path and returns exact
\texttt{Sqrt[\ldots]} forms. \B{SchurDecomposition} is numeric only---a generic
matrix has no closed-form symbolic Schur form---and is left unevaluated on symbolic
input rather than answered wrongly (\texttt{src/linalg/ludecomp.c},
\texttt{qrdecomp.c}, \texttt{svdecomp.c}, \texttt{schurdecomp.c},
\texttt{lapack.c}).}

\subsection{Eigenvalues, eigenvectors, and the Jordan form}
\label{ssec:la-eigen}

An eigenvector is a direction the matrix merely scales; its eigenvalue is the
scale. Finding them is the deepest structural question you can ask of a linear map,
and it is where the exact/numeric split is most vivid.

\begin{codepairs}
\pair{linear-algebra/eigen/1}{\B{CharacteristicPolynomial} gives \(\det(A - xI)\),
whose roots are the eigenvalues: \(x^2 - 5x - 2\) here.}
\pair{linear-algebra/eigen/2}{\B{Eigenvalues} of an exact matrix stays exact,
returning genuine surds---and a \(0\), certifying the matrix is singular.}
\pair{linear-algebra/eigen/3}{Symbolically, the \(2\times2\) eigenvalues are the
trace-and-determinant quadratic formula.}
\pair{linear-algebra/eigen/4}{A \(90^\circ\) rotation has no real eigenvectors, so
its eigenvalues are the complex pair \(\pm i\).}
\pair{linear-algebra/eigen/5}{\B{Eigenvectors} of a defective matrix: the repeated
eigenvalue \(2\) supplies only one eigenvector, so the deficient slot is padded
with a zero vector to keep the list aligned with the eigenvalues.}
\end{codepairs}

Pair~5 is the crack through which the Jordan form enters. The matrix there has
eigenvalue \(2\) with algebraic multiplicity \(2\) but only \emph{one} independent
eigenvector---it is \emph{defective}, not diagonalizable. The zero-vector padding
is Mathilda telling you honestly that an eigenvector is missing. The remedy is the
\emph{Jordan canonical form}\index{Jordan canonical form}: the closest a defective
matrix can come to diagonal, with the missing eigenvectors replaced by
\emph{generalized} ones\index{generalized eigenvector} and a tell-tale \(1\) on the
superdiagonal.

\begin{codepairs}
\pair{linear-algebra/eigen/6}{A matrix with a defective eigenvalue \(12\).}
\pair{linear-algebra/eigen/7}{\B{JordanDecomposition} returns the similarity
\(S\) and the Jordan form \(J\)\ldots}
\pair{linear-algebra/eigen/8}{\ldots whose \(J\) has a \(2\times2\) Jordan block for
\(12\): the eigenvalue twice on the diagonal, a \(1\) above it marking the deficient
eigenspace.}
\pair{linear-algebra/eigen/9}{And the decomposition reconstructs the matrix:
\(A = S\,J\,S^{-1}\).}
\end{codepairs}

The superdiagonal \(1\) in pair~8 is the entire content of the Jordan form: it says
that eigenvalue \(12\) has a two-step chain of generalized eigenvectors rather than
two independent ones. Read \(J\) and you read the matrix's deepest invariant---the
sizes of its Jordan blocks---which no similarity transformation can change.

\calloutnote{Under the hood}{\B{Eigenvalues} does \emph{not} form \(\det(A-xI)\) and
expand a symbolic determinant; it uses the Faddeev--LeVerrier--Souriau
recurrence\index{Faddeev-LeVerrier algorithm}, which produces the characteristic
polynomial in \(O(n^4)\) matrix multiplications, then hands the polynomial to
\B{Solve}---so eigenvalues of a large exact matrix return quickly, and the exact
roots come back as radicals or \B{Root} objects (\S\ref{sec:algebra}). The exact
\B{JordanDecomposition} reads each block's size from the nullity sequence
\(\dim\ker(A-\lambda I)^k\); a machine-real matrix with a distinct spectrum is
diagonalizable and takes a direct LAPACK \texttt{dgeev} path (a \(100\times100\)
in \(\approx\!9.6\)~ms), while a defective or ambiguous numeric matrix is
rationalised, decomposed exactly, and numericalised back to recover the true block
structure (\texttt{src/linalg/eigen.c}, \texttt{charpoly.c}, \texttt{jordandecomp.c}).}

\usagebox{Eigenvalues}

\usagebox{JordanDecomposition}

\subsection{The machine--exact divide}
\label{ssec:la-precision}

Every routine in this section has, in fact, been choosing between worlds all along.
It is worth watching the choice happen on one matrix. The Hilbert matrix---entries
\(1/(i+j-1)\), symmetric and famously ill-conditioned---is the perfect specimen.

\begin{codepairs}
\pair{linear-algebra/precision/1}{Exact: \B{Det} of the \(5\times5\)
\B{HilbertMatrix} is the precise rational \(1/266716800000\).}
\pair{linear-algebra/precision/2}{Machine: the same determinant over doubles is a
single \texttt{Real}---accurate to the digits shown, but a rounded number, not an
exact one.}
\pair{linear-algebra/precision/3}{Arbitrary precision: at \(40\) digits (MPFR) the
determinant of the \(7\times7\) Hilbert matrix comes back to full requested
precision.}
\pair{linear-algebra/precision/4}{Exact inversion of an ill-conditioned matrix is
no trouble at all when the arithmetic is exact---integer entries, computed by
FLINT.}
\pair{linear-algebra/precision/5}{The visible packed form: an \B{NDArray} is a
dense machine buffer that announces itself.}
\pair{linear-algebra/precision/6}{\B{Dot} over two \texttt{NDArray}s runs on the raw
buffers and returns a closed-system \texttt{NDArray}\ldots}
\pair{linear-algebra/precision/7}{\ldots and \B{Det} of one returns a bare machine
scalar, straight off the buffer.}
\end{codepairs}

The three determinants---an exact rational, a machine \texttt{Real}, and a
\(40\)-digit MPFR number---are the same mathematical quantity computed in three
number systems, and Mathilda selects among them purely from the \emph{kind} of the
entries. Exact in, exact out (via FLINT or fraction-free elimination); machine in,
LAPACK out; high-precision MPFR in, an MPFR kernel out. Nothing about the syntax
changes; the representation of the data decides the path.\index{arbitrary precision}\index{LAPACK}

\calloutnote{Performance}{The packed array is what makes the machine path fast, and
the rule that keeps it safe is the \emph{transparency gate}\index{packed array!transparency gate}: a
machine-number list may be held as a raw buffer, and a numeric routine on the
``aware'' list (\B{Dot}, \B{Det}, \B{Inverse}, \B{LinearSolve}, \B{Norm},
\B{MatrixRank}, and their kin) reads that buffer directly---no per-element boxing,
no re-parsing---so a \(300\times300\) \texttt{MatrixPower[A, 4]} is the \(0.69\)~ms of
its two \texttt{dgemm} calls. A routine \emph{not} on that list transparently
un-packs first, trading speed for correctness but never the answer. The same
machine heads lower inside \B{Compile}, so a linear-algebra kernel written at the
prompt runs at full speed inside a compiled numerical loop
(\texttt{src/pack.c}, \texttt{src/linalg/ndlinalg.c}; see also
\S\ref{sec:statistics} on packed reductions).}

\subsection*{Where this connects}

Linear algebra is the substrate the rest of the system computes on. The
least-squares solve of \S\ref{ssec:la-linsolve} is the machinery behind \B{Fit}
and the correlation matrix of \S\ref{sec:statistics}, both of them a \B{Dot} of a
matrix with its \B{Transpose}. The exact eigenvalues of \S\ref{ssec:la-eigen} flow
straight out of the polynomial solving and \B{Root} objects of \S\ref{sec:algebra},
and the fraction-free elimination beneath \B{RowReduce} is the same idea as the
subresultant sequences that tame coefficient growth there. The machine kernels and
the packed substrate reappear wherever numbers are crunched in bulk---the numerical
calculus, the optimizers, the machine-learning primitives---and every machine head
named here lowers inside \B{Compile} (Chapter~\ref{ch:compilation}), so the
linear algebra you write interactively is the linear algebra that runs in a tight
loop. For the exhaustive options of any function above, its reference page is one
\texttt{?Name} away at the prompt.
